\documentclass[10pt,a4paper]{amsart}
\usepackage{amsmath,amssymb,amsthm,graphicx}
\usepackage{verbatim}
\newenvironment{reference}{\comment}{\endcomment}
%\newenvironment{reference}{}{}
\newenvironment{slogan}{\comment}{\endcomment}
\newenvironment{history}{\comment}{\endcomment}

% For commutative diagrams we use Xy-pic
\usepackage[all]{xy}

% We use 2cell for 2-commutative diagrams.
\xyoption{2cell}
% Theorem environments.
%
\theoremstyle{plain}
\newtheorem{theorem}[subsection]{Theorem}
\newtheorem{proposition}[subsection]{Proposition}
\newtheorem{lemma}[subsection]{Lemma}

\theoremstyle{definition}
\newtheorem{definition}[subsection]{Definition}
\newtheorem{example}[subsection]{Example}
\newtheorem{exercise}[subsection]{Exercise}
\newtheorem{situation}[subsection]{Situation}

\theoremstyle{remark}
\newtheorem{remark}[subsection]{Remark}
\newtheorem{remarks}[subsection]{Remarks}

\numberwithin{equation}{subsection}

% Macros
%
\def\lim{\mathop{\mathrm{lim}}\nolimits}
\def\colim{\mathop{\mathrm{colim}}\nolimits}
\def\Spec{\mathop{\mathrm{Spec}}}
\def\Hom{\mathop{\mathrm{Hom}}\nolimits}
\def\Ext{\mathop{\mathrm{Ext}}\nolimits}
\def\SheafHom{\mathop{\mathcal{H}\!\mathit{om}}\nolimits}
\def\SheafExt{\mathop{\mathcal{E}\!\mathit{xt}}\nolimits}
\def\Sch{\mathit{Sch}}
\def\Mor{\mathop{\mathrm{Mor}}\nolimits}
\def\Ob{\mathop{\mathrm{Ob}}\nolimits}
\def\Sh{\mathop{\mathit{Sh}}\nolimits}
\def\NL{\mathop{N\!L}\nolimits}
\def\CH{\mathop{\mathrm{CH}}\nolimits}
\def\proetale{{pro\text{-}\acute{e}tale}}
\def\etale{{\acute{e}tale}}
\def\QCoh{\mathit{QCoh}}
\def\Ker{\mathop{\mathrm{Ker}}}
\def\Im{\mathop{\mathrm{Im}}}
\def\Coker{\mathop{\mathrm{Coker}}}
\def\Coim{\mathop{\mathrm{Coim}}}

% Boxtimes
%
\DeclareMathSymbol{\boxtimes}{\mathbin}{AMSa}{"02}

%
% Macros for moduli stacks/spaces
%
\def\QCohstack{\mathcal{QC}\!\mathit{oh}}
\def\Cohstack{\mathcal{C}\!\mathit{oh}}
\def\Spacesstack{\mathcal{S}\!\mathit{paces}}
\def\Quotfunctor{\mathrm{Quot}}
\def\Hilbfunctor{\mathrm{Hilb}}
\def\Curvesstack{\mathcal{C}\!\mathit{urves}}
\def\Polarizedstack{\mathcal{P}\!\mathit{olarized}}
\def\Complexesstack{\mathcal{C}\!\mathit{omplexes}}
% \Pic is the operator that assigns to X its picard group, usage \Pic(X)
% \Picardstack_{X/B} denotes the Picard stack of X over B
% \Picardfunctor_{X/B} denotes the Picard functor of X over B
\def\Pic{\mathop{\mathrm{Pic}}\nolimits}
\def\Picardstack{\mathcal{P}\!\mathit{ic}}
\def\Picardfunctor{\mathrm{Pic}}
\def\Deformationcategory{\mathcal{D}\!\mathit{ef}}
\usepackage{cleveref}
\crefname{originalreference}{original reference}{original references}
\Crefname{originalreference}{Original reference}{Original references}
\makeatletter
\newlabel{lemma-flat-ascent-impurity}{{0CVA}{source}{}{sourcepointer-1}{}}
\newlabel{lemma-flat-ascent-impurity@cref}{{[originalreference][0][]0CVA}{[1][0][]source}{}{}{}}
\makeatother
\crefname{originalreference}{original reference}{original references}
\Crefname{originalreference}{Original reference}{Original references}
\makeatletter
\newlabel{lemma-impure-limit}{{05IZ}{source}{}{sourcepointer-2}{}}
\newlabel{lemma-impure-limit@cref}{{[originalreference][0][]05IZ}{[1][0][]source}{}{}{}}
\makeatother
\crefname{originallemma}{original lemma}{original lemmas}
\Crefname{originallemma}{Original lemma}{Original lemmas}
\makeatletter
\newlabel{more-morphisms-lemma-etale-makes-quasi-finite-finite-at-point}{{02LK}{source}{}{sourcepointer-3}{}}
\newlabel{more-morphisms-lemma-etale-makes-quasi-finite-finite-at-point@cref}{{[originallemma][0][]02LK}{[1][0][]source}{}{}{}}
\makeatother
\crefname{originallemma}{original lemma}{original lemmas}
\Crefname{originallemma}{Original lemma}{Original lemmas}
\makeatletter
\newlabel{morphisms-lemma-composition-quasi-finite}{{01TL}{source}{}{sourcepointer-4}{}}
\newlabel{morphisms-lemma-composition-quasi-finite@cref}{{[originallemma][0][]01TL}{[1][0][]source}{}{}{}}
\makeatother
\crefname{originalreference}{original reference}{original references}
\Crefname{originalreference}{Original reference}{Original references}
\makeatletter
\newlabel{situation-pre-pure}{{05FW}{source}{}{sourcepointer-5}{}}
\newlabel{situation-pre-pure@cref}{{[originalreference][0][]05FW}{[1][0][]source}{}{}{}}
\makeatother
\crefname{originallemma}{original lemma}{original lemmas}
\Crefname{originallemma}{Original lemma}{Original lemmas}
\makeatletter
\newlabel{spaces-divisors-lemma-base-change-relative-assassin}{{0CV2}{source}{}{sourcepointer-6}{}}
\newlabel{spaces-divisors-lemma-base-change-relative-assassin@cref}{{[originallemma][0][]0CV2}{[1][0][]source}{}{}{}}
\makeatother
\crefname{originallemma}{original lemma}{original lemmas}
\Crefname{originallemma}{Original lemma}{Original lemmas}
\makeatletter
\newlabel{spaces-morphisms-lemma-base-change-quasi-finite}{{03XL}{source}{}{sourcepointer-7}{}}
\newlabel{spaces-morphisms-lemma-base-change-quasi-finite@cref}{{[originallemma][0][]03XL}{[1][0][]source}{}{}{}}
\makeatother
\crefname{originallemma}{original lemma}{original lemmas}
\Crefname{originallemma}{Original lemma}{Original lemmas}
\makeatletter
\newlabel{spaces-morphisms-lemma-etale-locally-quasi-finite}{{03XU}{source}{}{sourcepointer-8}{}}
\newlabel{spaces-morphisms-lemma-etale-locally-quasi-finite@cref}{{[originallemma][0][]03XU}{[1][0][]source}{}{}{}}
\makeatother
\crefname{originallemma}{original lemma}{original lemmas}
\Crefname{originallemma}{Original lemma}{Original lemmas}
\makeatletter
\newlabel{spaces-properties-lemma-cofinal-etale}{{04JW}{source}{}{sourcepointer-9}{}}
\newlabel{spaces-properties-lemma-cofinal-etale@cref}{{[originallemma][0][]04JW}{[1][0][]source}{}{}{}}
\makeatother
\title{Purity along a geometric point: four equivalent formulations}
\author{The Stacks Project Authors}
\date{stacks-project:spaces-flat.tex:local-snapshot-20261012; GNU FDL 1.2 or later}
\begin{document}
\maketitle
\noindent\textit{Editorial note.} All delivered mathematical text is the author's own, copied from the pinned Stacks Project chapter file spaces-flat.tex; no statement or proof step was rewritten, reordered or removed.\par\medskip
\noindent\textit{Editorial note.} The counted proof is the author's complete argument: a limit argument descending purity from the etale local ring to an etale neighbourhood, and an ascent along a locally quasi-finite base change using that a quasi-finite morphism over a strictly henselian base becomes finite\par\medskip
\noindent\textit{Editorial note.} The label lemma-pure-along-X-y names the counted statement, and the proof environment immediately following it in the source is delivered verbatim as its proof.\par\medskip
\noindent\textit{Editorial note.} The counted statement is carried by the author's own lemma environment in the chapter source.\par\medskip
\noindent\textit{Allowed standard prerequisites (editorial).}\begin{itemize}
\item Impurities as witnesses that a module is not pure along a point, and the definition of purity along a point in terms of the associated points of the geometric fibres
\item The etale local ring at a geometric point as the filtered colimit of the rings of etale neighbourhoods, and the limit characterisation of purity that follows from it
\item Ascent of an impurity along a locally quasi-finite base change, and the fact that a quasi-finite morphism over a strictly henselian base becomes finite after passing to an open and closed neighbourhood
\item Base change of the relative assassin under a finite morphism, and closedness of the induced map on underlying spaces
\end{itemize}\medskip
\ifdefined\hypertarget\hypertarget{sourcepointer-1}{}\fi
\noindent\textit{Source reference.} Original reference 0CVA: sources/stacks-spaces-flat/pinned-original-source.tex, source lines 366--366. Prerequisite pointer invoked at proof line(s) 444; it names a standard fact of the same chapter that the counted argument uses, and every step of the counted proof built on it is retained verbatim in the delivered body.\par\medskip
\ifdefined\hypertarget\hypertarget{sourcepointer-2}{}\fi
\noindent\textit{Source reference.} Original reference 05IZ: sources/stacks-spaces-flat/pinned-original-source.tex, source lines 319--319. Prerequisite pointer invoked at proof line(s) 433; it names a standard fact of the same chapter that the counted argument uses, and every step of the counted proof built on it is retained verbatim in the delivered body.\par\medskip
\ifdefined\hypertarget\hypertarget{sourcepointer-3}{}\fi
\noindent\textit{Source reference.} Original lemma 02LK: sources/stacks-chow/tags.txt, source lines 3411--3411. Prerequisite pointer invoked at proof line(s) 454; it names a standard fact of a sibling chapter that the counted argument uses, and every step of the counted proof built on it is retained verbatim in the delivered body.\par\medskip
\ifdefined\hypertarget\hypertarget{sourcepointer-4}{}\fi
\noindent\textit{Source reference.} Original lemma 01TL: sources/stacks-chow/tags.txt, source lines 2378--2378. Prerequisite pointer invoked at proof line(s) 451; it names a standard fact of a sibling chapter that the counted argument uses, and every step of the counted proof built on it is retained verbatim in the delivered body.\par\medskip
\ifdefined\hypertarget\hypertarget{sourcepointer-5}{}\fi
\noindent\textit{Source reference.} Original reference 05FW: sources/stacks-spaces-flat/pinned-original-source.tex, source lines 260--260. Prerequisite pointer invoked at proof line(s) 403; it names a standard fact of the same chapter that the counted argument uses, and every step of the counted proof built on it is retained verbatim in the delivered body.\par\medskip
\ifdefined\hypertarget\hypertarget{sourcepointer-6}{}\fi
\noindent\textit{Source reference.} Original lemma 0CV2: sources/stacks-chow/tags.txt, source lines 16122--16122. Prerequisite pointer invoked at proof line(s) 460; it names a standard fact of a sibling chapter that the counted argument uses, and every step of the counted proof built on it is retained verbatim in the delivered body.\par\medskip
\ifdefined\hypertarget\hypertarget{sourcepointer-7}{}\fi
\noindent\textit{Source reference.} Original lemma 03XL: sources/stacks-chow/tags.txt, source lines 5152--5152. Prerequisite pointer invoked at proof line(s) 450; it names a standard fact of a sibling chapter that the counted argument uses, and every step of the counted proof built on it is retained verbatim in the delivered body.\par\medskip
\ifdefined\hypertarget\hypertarget{sourcepointer-8}{}\fi
\noindent\textit{Source reference.} Original lemma 03XU: sources/stacks-chow/tags.txt, source lines 5161--5161. Prerequisite pointer invoked at proof line(s) 438; it names a standard fact of a sibling chapter that the counted argument uses, and every step of the counted proof built on it is retained verbatim in the delivered body.\par\medskip
\ifdefined\hypertarget\hypertarget{sourcepointer-9}{}\fi
\noindent\textit{Source reference.} Original lemma 04JW: sources/stacks-chow/tags.txt, source lines 5952--5952. Prerequisite pointer invoked at proof line(s) 430; it names a standard fact of a sibling chapter that the counted argument uses, and every step of the counted proof built on it is retained verbatim in the delivered body.\par\medskip
\begin{lemma}
\label{lemma-pure-along-X-y}
In Situation \ref{situation-pre-pure}. Let $\overline{y}$ be a geometric
point lying over $y$. Let $\mathcal{O} = \mathcal{O}_{Y, \overline{y}}$
be the \'etale local ring of $Y$ at $\overline{y}$. Denote
$Y^{sh} = \Spec(\mathcal{O})$, $X^{sh} = X \times_Y Y^{sh}$, and
$\mathcal{F}^{sh}$ the pullback of $\mathcal{F}$ to $X^{sh}$.
The following are equivalent
\begin{enumerate}
\item there exists an impurity
$(Y^{sh} \to Y, y' \leadsto \overline{y}, \xi)$
of $\mathcal{F}$ above $y$,
\item every point of $\text{Ass}_{X^{sh}/Y^{sh}}(\mathcal{F}^{sh})$
specializes to a point of the closed fibre $X_{\overline{y}}$,
\item there exists an impurity $(T \to Y, t' \leadsto t, \xi)$
of $\mathcal{F}$ above $y$ such that $(T, t) \to (Y, y)$ is an
\'etale neighbourhood, and
\item there exists an impurity $(T \to Y, t' \leadsto t, \xi)$
of $\mathcal{F}$ above $y$ such that $T \to Y$ is quasi-finite at $t$.
\end{enumerate}
\end{lemma}

\begin{proof}
That parts (1) and (2) are equivalent is immediate from the definition.

\medskip\noindent
Recall that $\mathcal{O} = \mathcal{O}_{Y, \overline{y}}$
is the filtered colimit of $\mathcal{O}(V)$ over the category
of \'etale neighbourhoods $(V, \overline{v}) \to (Y, \overline{y})$
(Properties of Spaces, Lemma \ref{spaces-properties-lemma-cofinal-etale}).
Moreover, it suffices to consider affine \'etale neighbourhoods $V$.
Hence $Y^{sh} = \Spec(\mathcal{O}) = \lim \Spec(\mathcal{O}(V)) = \lim V$.
Thus we see that (1) implies (3) by Lemma \ref{lemma-impure-limit}.

\medskip\noindent
Since an \'etale morphism is locally quasi-finite
(Morphisms of Spaces, Lemma
\ref{spaces-morphisms-lemma-etale-locally-quasi-finite})
we see that (3) implies (4).

\medskip\noindent
Finally, assume (4). After replacing $T$ by an open neighbourhood of $t$
we may assume $T \to Y$ is locally quasi-finite.
By Lemma \ref{lemma-flat-ascent-impurity}
we find an impurity
$(T_1 \to Y^{sh}, t_1' \leadsto t_1, \xi_1)$
with $T_1 \to T \times_Y Y^{sh}$
\'etale. Since an \'etale morphism is locally quasi-finite
and using Morphisms of Spaces, Lemma
\ref{spaces-morphisms-lemma-base-change-quasi-finite} and
Morphisms, Lemma \ref{morphisms-lemma-composition-quasi-finite}
we see that $T_1 \to Y^{sh}$ is locally quasi-finite.
As $\mathcal{O}$ is strictly henselian, we can apply More on Morphisms, Lemma
\ref{more-morphisms-lemma-etale-makes-quasi-finite-finite-at-point}
to see that after replacing $T_1$ by an open and closed neighbourhood
of $t_1$ we may assume that $T_1 \to Y^{sh} = \Spec(\mathcal{O})$
is finite. Let $\theta \in |X^{sh}|$ be the image of
$\xi_1$ and let $y' \in \Spec(\mathcal{O})$ be the image
of $t_1'$. By Divisors on Spaces, Lemma
\ref{spaces-divisors-lemma-base-change-relative-assassin}
we see that $\theta \in \text{Ass}_{X^{sh}/Y^{sh}}(\mathcal{F}^{sh})$.
Since $\pi : X_{T_1} \to X^{sh}$
is finite, it induces a closed map $|X_{T_1}| \to |X^{sh}|$.
Hence the image of $\overline{\{\xi_1\}}$ is $\overline{\{\theta\}}$.
It follows that $(Y^{sh} \to Y, y' \leadsto \overline{y}, \theta)$
is an impurity of $\mathcal{F}$ above $y$ and the proof is complete.
\end{proof}

\end{document}
